%% Formal sub-subleading Virasoro recurrence and its geometric comparison.

\providecommand{\Vir}{\mathrm{Vir}}
\providecommand{\ClaimStatusProvedHere}{}
\providecommand{\ClaimStatusConditional}{}
\providecommand{\ClaimStatusDefinitional}{}
\providecommand{\ClaimStatusOpen}{}

\chapter{The Sub-subleading Virasoro Recurrence}
\label{ch:shadow-tower-sub-subleading-platonic}

\section*{Formal sequence and residue sequence}

This chapter computes one coefficient of a formal quadratic
recurrence.  Its geometric interpretation belongs to a separate
comparison package.  Let \(R_r(c)\) denote the formal sequence and let
\(S_r(\Vir_c;\mathsf H_{\mathrm{res}})\) denote the scalar extracted
from an ordered residue package.  The package
\(\mathsf H_{\mathrm{asymp}}\) consists of an identification
\(S_r=R_r\) together with a common large-\(c\) domain and a termwise
Laurent expansion.  Every asymptotic assertion about the geometric
residue tower below carries
\(\mathsf H_{\mathrm{res}}+\mathsf H_{\mathrm{asymp}}\).

\section{The formal Laurent problem}
\label{sec:stss-setup}

\begin{definition}[Rescaled formal sequence]
\label{def:phi-laurent}
\ClaimStatusDefinitional
Set \(R_3(c)=2\), \(R_4(c)=10/[c(5c+22)]\), and for \(r\ge5\) define
\[
 R_r(c)=-\frac1{rc}
 \sum_{\substack{j+k=r+2\\3\le j\le k<r}}
 \epsilon_{j,k}\,jkR_j(c)R_k(c),
 \qquad
 \epsilon_{j,k}=\begin{cases}1,&j<k,\\[2pt]\frac12,&j=k.\end{cases}
\]
Write \(\Phi_r(c)=c^{r-2}R_r(c)\).  Its formal Laurent expansion is
\[
 \Phi_r(c)=A_r+\frac{B_r}{c}+\frac{\Gamma_r}{c^2}+O(c^{-3}).
\]
\end{definition}

\begin{lemma}[Rescaled recurrence]
\label{lem:phi-recurrence}
\ClaimStatusProvedHere
For \(r\ge5\),
\[
 \Phi_r=-\frac{6(r-1)}r\Phi_{r-1}
 -\frac1{rc}
 \sum_{\substack{j+k=r+2\\4\le j\le k\le r-2}}
 \epsilon_{j,k}\,jk\Phi_j\Phi_k.
\]
\end{lemma}

\begin{proof}
Substitute \(R_j=c^{2-j}\Phi_j\) into the defining recurrence and
separate the pair \((3,r-1)\).
\end{proof}

The leading two formal coefficients are
\[
 A_r=\frac{8(-6)^{r-4}}r,
 \qquad
 B_r=-A_r\left(\frac{22}{5}
 +\frac{(r-4)(r-5)}{18}\right).
\]
They may also be taken as the hypotheses for the following coefficient
extraction.

\section{The coefficient recurrence}
\label{sec:stss-recurrence}

\begin{proposition}[Formal recurrence for \(\Gamma_r\)]
\label{prop:gamma-recurrence}
\ClaimStatusProvedHere
For \(r\ge5\),
\[
 \Gamma_r=-\frac{6(r-1)}r\Gamma_{r-1}-\sigma_r^{(\Gamma)},
\]
where
\[
 \sigma_r^{(\Gamma)}=
 \frac1r\sum_{\substack{j+k=r+2\\4\le j\le k\le r-2}}
 \epsilon_{j,k}\,jk(A_jB_k+B_jA_k).
\]
\end{proposition}

\begin{proof}
Extract the coefficient of \(c^{-2}\) in
Lemma~\ref{lem:phi-recurrence}.  The external factor \(c^{-1}\)
selects the coefficient \(A_jB_k+B_jA_k\) from each product.
\end{proof}

\begin{lemma}[Closed source ratio]
\label{lem:gamma-source-ratio-closed-form}
\ClaimStatusProvedHere
For \(r\ge5\),
\[
 \frac{\sigma_r^{(\Gamma)}}{A_r}
 =-\frac{44(r-5)}{45}
 -\frac{(r-5)(r-6)(r-7)}{243}.
\]
\end{lemma}

\begin{proof}
The identity \(jA_j=8(-6)^{j-4}\) makes
\(jkA_jA_k=64(-6)^{r-6}\) independent of the partition
\(j+k=r+2\).  Symmetrization converts the polynomial part into
\[
 \sum_{j=4}^{r-2}(j-4)(j-5)
 =\frac{(r-5)(r-6)(r-7)}3.
\]
Substitution of \(B_j/A_j\) gives the displayed ratio.
\end{proof}

\section{Closed form and geometric comparison}
\label{sec:stss-closed-form}

\begin{theorem}[Formal sub-subleading coefficient]
\label{thm:shadow-tower-sub-subleading-closed-form}
\ClaimStatusProvedHere
For \(r\ge4\),
\[
 \frac{\Gamma_r}{A_r}
 =\frac{484}{25}
 +\frac{22(r-4)(r-5)}{45}
 +\frac{(r-4)(r-5)(r-6)(r-7)}{972}.
\]
Equivalently, with
\(\phi_r=22/5+(r-4)(r-5)/18\),
\[
 \frac{\Gamma_r}{A_r}
 =\phi_r^2-
 \frac{(r-4)(r-5)(r^2-7r+9)}{486}.
\]
\end{theorem}

\begin{proof}
Since \(A_r=-6(r-1)A_{r-1}/r\), division by \(A_r\) turns
Proposition~\ref{prop:gamma-recurrence} into a first-difference
equation.  Starting from \(\Gamma_4/A_4=484/25\), use
Lemma~\ref{lem:gamma-source-ratio-closed-form} and the identities
\[
 \sum_{s=5}^{r}(s-5)=\frac{(r-4)(r-5)}2,
 \quad
 \sum_{s=5}^{r}(s-5)(s-6)(s-7)
 =\frac{(r-4)(r-5)(r-6)(r-7)}4.
\]
The second formula follows by expansion.
\end{proof}

\begin{remark}[Initial value]
\label{rem:gamma-base-case}
At \(r=4,5\) the polynomial increments vanish and
\(\Gamma_r/A_r=484/25\).  This is the square of the initial
subleading ratio.
\end{remark}

\begin{remark}[Relation with a Riccati model]
\label{rem:gamma-riccati-origin}
The formula is intrinsic to the recurrence of
Definition~\ref{def:phi-laurent}.  A Riccati generating function gives
another presentation of the same formal sequence once its coefficient
comparison has been proved.
\end{remark}

\begin{remark}[Residue-tower interpretation]
\ClaimStatusConditional
Under \(\mathsf H_{\mathrm{res}}+\mathsf H_{\mathrm{asymp}}\), the
coefficient \(\Gamma_r\) is the sub-subleading Laurent coefficient of
\(c^{r-2}S_r(\Vir_c;\mathsf H_{\mathrm{res}})\).  Constructing this
package is the geometric proof obligation.
\end{remark}

\section{Congruence data of the formal quartic}
\label{sec:stss-kummer}

\begin{lemma}[Integral numerator]
\label{lem:gamma-numerator-quartic-polynomial}
\ClaimStatusProvedHere
Writing \(\Gamma_r/A_r=N(r)/24300\) gives
\[
 N(r)=25r^4-550r^3+16355r^2-122870r+729048.
\]
\end{lemma}

\begin{proof}
Put the three terms of
Theorem~\ref{thm:shadow-tower-sub-subleading-closed-form} over the
common denominator \(24300\) and expand.
\end{proof}

\begin{lemma}[Irreducibility of the numerator]
\label{lem:gamma-numerator-irreducible}
\ClaimStatusProvedHere
The polynomial \(N\) is primitive and irreducible in
\(\mathbb Q[r]\).
\end{lemma}

\begin{proof}
Its reduction modulo \(7\) is
\[
 4r^4+3r^3+3r^2+r+5\in\mathbb F_7[r].
\]
The finite-field irreducibility criterion gives
\[
 \gcd\!\left(N(r),r^{49}-r\right)=1,
 \qquad
 r^{2401}-r\equiv0\pmod{N(r),7}.
\]
Thus its reduction has degree four over \(\mathbb F_7\).  Gauss's
lemma lifts irreducibility to \(\mathbb Q\).
\end{proof}

\begin{remark}[The residue class modulo \(691\)]
\label{rem:gamma-691-emergence-sporadic}
\ClaimStatusProvedHere
The exact factorization
\[
 N(8)=613608=2^3\cdot3\cdot37\cdot691
\]
is a congruence of the formal recurrence.  The roots of \(N\) in
\(\mathbb F_{691}\) are
\(8,315,423,658\).  A Bernoulli or Kummer interpretation requires a
separate comparison package \(\mathsf H_{\mathrm{Kum}}\) connecting
this recurrence to a Bernoulli-valued geometric invariant.
\end{remark}

\begin{remark}[Prime values of an integral quartic]
\label{rem:gamma-irregular-primes-dense-but-structureless}
The prime divisors of \(N(r)\) record the arithmetic of this integral
quartic.  Their interpretation as geometric Kummer data begins after
\(\mathsf H_{\mathrm{Kum}}\) supplies the relevant motive, integral
lattice, and reduction map.
\end{remark}

\section{Verification surface}
\label{sec:stss-verification}

The formal identities admit four independent checks: coefficient
extraction from the recurrence, substitution of the closed form,
symbolic expansion of \(N(r)\), and finite-field factorization.  The
geometric lane has two additional checks: construction of
\(\mathsf H_{\mathrm{res}}\) and uniform control of the Laurent
expansion in \(\mathsf H_{\mathrm{asymp}}\).

\section{Status}
\label{sec:stss-summary}

The quartic formula is proved for the formal quadratic recurrence.
Its Virasoro residue interpretation is conditional on
\(\mathsf H_{\mathrm{res}}+\mathsf H_{\mathrm{asymp}}\), and its
Bernoulli interpretation is governed by \(\mathsf H_{\mathrm{Kum}}\).
